%=============================================================================
% Chapter 1: Functions in C++
% Structured Programming with C++ -- Lectures 11-12
%=============================================================================

\chapter{Functions in C++}
\label{ch:functions}

%-----------------------------------------------------------------------------
\section{Introduction to Functions}
\label{sec:intro-functions}
%-----------------------------------------------------------------------------

\begin{definitionbox}{Function}
A \textbf{function} is a self-contained set of statements designed to accomplish
a particular task. Every C++ program contains at least one function ---
\texttt{main()} --- and may define any number of additional functions.
\end{definitionbox}

In practice, large programs are difficult to write, read, and maintain as a
single monolithic block of code. The solution is \emph{modular programming}:
dividing a program into smaller, manageable pieces called \textbf{functions}.
Each function performs a specific, well-defined task, and the complete program
is built by combining these functions together.

\begin{infobox}{Advantages of Using Functions}
\begin{enumerate}
    \item \textbf{Easier to write:} Each function focuses on a single task,
          making the logic simpler to implement.
    \item \textbf{Easier to read:} Well-named functions make the overall
          program flow self-documenting and understandable at a glance.
    \item \textbf{Easier to debug:} Errors can be isolated to individual
          functions, reducing the search space during troubleshooting.
    \item \textbf{Reusability:} A function can be called multiple times with
          different parameters, eliminating code duplication.
    \item \textbf{Team development:} Different programmers can work on
          different functions simultaneously.
\end{enumerate}
\end{infobox}

\begin{conceptbox}{Modular Programming}
Modular programming is a software design technique that decomposes a program
into separate sub-programs (modules or functions). Each module is developed
and tested independently, then integrated into the complete program. This
approach reduces complexity and improves code quality.
\end{conceptbox}

%-----------------------------------------------------------------------------
\section{Defining a Function}
\label{sec:defining-function}
%-----------------------------------------------------------------------------

Every function definition in C++ consists of four parts: a \textbf{return type},
a \textbf{function name}, a \textbf{parameter list}, and a \textbf{function
body}. The general form is:

\begin{cppcode}
return-type function-name(parameter-list)
{
    // body of the function
    statement1;
    statement2;
    ...
    return value;  // (if return-type is not void)
}
\end{cppcode}

\begin{importantbox}{Components of a Function Definition}
\begin{itemize}
    \item \textbf{Return type:} The data type of the value returned by the
          function. Common types include \texttt{int}, \texttt{float},
          \texttt{double}, \texttt{char}, and \texttt{void} (when the function
          returns nothing).
    \item \textbf{Function name:} A valid C++ identifier that describes the
          purpose of the function.
    \item \textbf{Parameter list:} A comma-separated list of typed variables
          that receive values from the caller. May be empty if the function
          requires no input.
    \item \textbf{Function body:} The block of statements enclosed in braces
          \texttt{\{\}} that define what the function does.
\end{itemize}
\end{importantbox}

%-----------------------------------------------------------------------------
\subsection{Function Prototype (Declaration)}
\label{subsec:function-prototype}
%-----------------------------------------------------------------------------

A \textbf{function prototype} tells the compiler about a function's name,
return type, and parameters \emph{before} the function is fully defined. This
allows \texttt{main()} (or any calling function) to appear before the function
definition in the source file.

\begin{definitionbox}{Function Prototype}
A \textbf{function prototype} is a declaration of a function that specifies its
return type, name, and parameter types, terminated by a semicolon.
It does \emph{not} include the function body. The general form is:

\medskip
\texttt{return-type function-name(parameter-type-list);}
\end{definitionbox}

\begin{cppcode}
#include <iostream>
using namespace std;

// Function prototype (declaration)
int add(int, int);

int main()
{
    int result = add(10, 20);
    cout << "Sum = " << result << endl;
    return 0;
}

// Function definition (after main)
int add(int a, int b)
{
    return a + b;
}
\end{cppcode}

\noindent\textbf{Output:}
\begin{verbatim}
Sum = 30
\end{verbatim}

\begin{warningbox}{Function Declaration Order}
In C++, a function must be \emph{declared} before it is called. You can either
place the full function definition before \texttt{main()}, or provide a
\textbf{function prototype} (declaration) before \texttt{main()} and define
the function afterwards. If you call a function without a prior declaration,
the compiler will report an error.
\end{warningbox}

%-----------------------------------------------------------------------------
\subsection{Function Return Types Summary}
%-----------------------------------------------------------------------------

The following table summarises the most common return types used with
functions:

\begin{center}
\renewcommand{\arraystretch}{1.4}
\begin{tabular}{@{}llp{6cm}@{}}
\toprule
\textbf{Return Type} & \textbf{Example Signature}
    & \textbf{Description} \\
\midrule
\texttt{void}   & \texttt{void greet()}
    & Performs an action; returns nothing. \\
\texttt{int}    & \texttt{int square(int n)}
    & Returns an integer result. \\
\texttt{float}  & \texttt{float aver(int a, int b)}
    & Returns a floating-point result. \\
\texttt{double} & \texttt{double area(double r)}
    & Returns a double-precision result. \\
\texttt{char}   & \texttt{char grade(int mark)}
    & Returns a single character. \\
\texttt{bool}   & \texttt{bool isEven(int n)}
    & Returns \texttt{true} or \texttt{false}. \\
\bottomrule
\end{tabular}
\end{center}

%-----------------------------------------------------------------------------
\subsection{Scope: Local vs.\ Global}
%-----------------------------------------------------------------------------

Variables declared inside a function are \textbf{local} to that function ---
they exist only while the function is executing and cannot be accessed from
outside. Variables declared outside all functions are \textbf{global} and can
be accessed from any function in the program. A detailed discussion of scope
follows in Chapter~\ref{ch:advanced-functions}.

%-----------------------------------------------------------------------------
\subsection{Example: A \texttt{void} Function}
%-----------------------------------------------------------------------------

The simplest kind of function performs an action without returning a value.
Its return type is \texttt{void}.

\begin{cppcode}
#include <iostream>
using namespace std;

// Function definition
void printMessage()
{
    cout << "Welcome to the world of C++ Functions!"
         << endl;
    cout << "Functions make programs modular and reusable."
         << endl;
}

int main()
{
    printMessage();   // Function call
    printMessage();   // Called again -- reusability in action
    return 0;
}
\end{cppcode}

\noindent\textbf{Output:}
\begin{verbatim}
Welcome to the world of C++ Functions!
Functions make programs modular and reusable.
Welcome to the world of C++ Functions!
Functions make programs modular and reusable.
\end{verbatim}

%-----------------------------------------------------------------------------
\subsection{Example: A Function that Returns a Value}
%-----------------------------------------------------------------------------

When a function computes and returns a result, its return type must match the
type of the value returned.

\begin{cppcode}
#include <iostream>
using namespace std;

// Function to find the maximum of two integers
int max(int a, int b)
{
    if (a > b)
        return a;
    else
        return b;
}

int main()
{
    int x = 15, y = 27;
    int result = max(x, y);
    cout << "The maximum of " << x << " and " << y
         << " is: " << result << endl;
    return 0;
}
\end{cppcode}

\noindent\textbf{Output:}
\begin{verbatim}
The maximum of 15 and 27 is: 27
\end{verbatim}

%-----------------------------------------------------------------------------
\section{The Return Statement}
\label{sec:return-statement}
%-----------------------------------------------------------------------------

\begin{definitionbox}{The \texttt{return} Statement}
The \texttt{return} statement terminates execution of the current function and
passes control back to the calling function. It has two forms:
\begin{itemize}
    \item \texttt{return;}\ --- used in \texttt{void} functions to exit early.
    \item \texttt{return expression;}\ --- used in non-\texttt{void} functions
          to send a value back to the caller.
\end{itemize}
\end{definitionbox}

\begin{conceptbox}{How \texttt{return} Works}
When the \texttt{return} statement executes, the function stops immediately.
Any statements after the \texttt{return} in the same block are \emph{not}
executed. In non-\texttt{void} functions, the returned expression is evaluated
and its value is made available at the point of the function call.
\end{conceptbox}

%-----------------------------------------------------------------------------
\subsection{Example: Calculating Squares of 1--10}
%-----------------------------------------------------------------------------

The following program uses a function to compute the square of an integer and
prints the squares of numbers from 1 to 10.

\begin{cppcode}
#include <iostream>
using namespace std;

int square(int y)
{
    return y * y;
}

int main()
{
    cout << "Number\tSquare" << endl;
    cout << "------\t------" << endl;

    for (int x = 1; x <= 10; x++)
    {
        cout << x << "\t" << square(x) << endl;
    }

    return 0;
}
\end{cppcode}

\noindent\textbf{Output:}
\begin{verbatim}
Number  Square
------  ------
1       1
2       4
3       9
4       16
5       25
6       36
7       49
8       64
9       81
10      100
\end{verbatim}

%-----------------------------------------------------------------------------
\subsection{Example: Calculating an Average}
%-----------------------------------------------------------------------------

This function receives two integer values and returns their average as a
floating-point number.

\begin{cppcode}
#include <iostream>
using namespace std;

float aver(int x1, int x2)
{
    float result;
    result = (x1 + x2) / 2.0;
    return result;
}

int main()
{
    int a, b;
    cout << "Enter two integers: ";
    cin >> a >> b;

    float avg = aver(a, b);
    cout << "The average of " << a << " and " << b
         << " is: " << avg << endl;

    return 0;
}
\end{cppcode}

\noindent\textbf{Sample Output:}
\begin{verbatim}
Enter two integers: 7 13
The average of 7 and 13 is: 10
\end{verbatim}

\begin{warningbox}{Integer Division Pitfall}
Note the use of \texttt{2.0} instead of \texttt{2} in the division. Dividing
two integers in C++ produces an integer result (truncating the decimal part).
Writing \texttt{2.0} forces floating-point division, ensuring an accurate
average.
\end{warningbox}

%-----------------------------------------------------------------------------
\subsection{Example: Sum of Squares Series}
%-----------------------------------------------------------------------------

The following program computes the sum of the series
$1^{2} + 2^{2} + 3^{2} + \cdots + n^{2}$ using a dedicated function.

\begin{cppcode}
#include <iostream>
using namespace std;

int summation(int x)
{
    int sum = 0;
    for (int i = 1; i <= x; i++)
    {
        sum += i * i;
    }
    return sum;
}

int main()
{
    int n;
    cout << "Enter a positive integer n: ";
    cin >> n;

    int result = summation(n);
    cout << "The sum of squares from 1 to " << n
         << " is: " << result << endl;

    return 0;
}
\end{cppcode}

\noindent\textbf{Sample Output:}
\begin{verbatim}
Enter a positive integer n: 5
The sum of squares from 1 to 5 is: 55
\end{verbatim}

\noindent The computation:
$1^{2} + 2^{2} + 3^{2} + 4^{2} + 5^{2} = 1 + 4 + 9 + 16 + 25 = 55$.

%-----------------------------------------------------------------------------
\subsection{Example: Maximum of Three Integers}
\label{subsec:max-of-three}
%-----------------------------------------------------------------------------

A function can call other functions. Here \texttt{max3()} finds the largest
of three integers by calling the two-argument \texttt{max()} function twice.

\begin{cppcode}
#include <iostream>
using namespace std;

int max(int a, int b)
{
    if (a > b)
        return a;
    else
        return b;
}

int max3(int a, int b, int c)
{
    return max(max(a, b), c);
}

int main()
{
    int x, y, z;
    cout << "Enter three integers: ";
    cin >> x >> y >> z;

    cout << "The largest value is: "
         << max3(x, y, z) << endl;

    return 0;
}
\end{cppcode}

\noindent\textbf{Sample Output:}
\begin{verbatim}
Enter three integers: 14 7 21
The largest value is: 21
\end{verbatim}

\begin{infobox}{Composing Functions}
Notice how \texttt{max3()} reuses \texttt{max()} rather than duplicating
comparison logic. This demonstrates the principle of \emph{function
composition}: building complex behaviour by combining simpler functions.
\end{infobox}

%-----------------------------------------------------------------------------
\subsection{Example: Logic Gates}
\label{subsec:logic-gates}
%-----------------------------------------------------------------------------

Functions can also model logical operations. The following program implements
the four basic logic gates --- AND, OR, XOR, and NOT --- as separate
functions.

\begin{cppcode}
#include <iostream>
using namespace std;

bool AND(bool a, bool b) { return a && b; }
bool OR(bool a, bool b)  { return a || b; }
bool XOR(bool a, bool b) { return a != b; }
bool NOT(bool a)          { return !a;     }

int main()
{
    bool x = true, y = false;

    cout << "x = " << x << ", y = " << y << endl;
    cout << "AND(x,y) = " << AND(x, y) << endl;
    cout << "OR(x,y)  = " << OR(x, y)  << endl;
    cout << "XOR(x,y) = " << XOR(x, y) << endl;
    cout << "NOT(x)   = " << NOT(x)    << endl;
    cout << "NOT(y)   = " << NOT(y)    << endl;

    return 0;
}
\end{cppcode}

\noindent\textbf{Output:}
\begin{verbatim}
x = 1, y = 0
AND(x,y) = 0
OR(x,y)  = 1
XOR(x,y) = 1
NOT(x)   = 0
NOT(y)   = 1
\end{verbatim}

\begin{conceptbox}{Logic Gate Truth Table}
\begin{center}
\renewcommand{\arraystretch}{1.3}
\begin{tabular}{@{}cccccc@{}}
\toprule
\textbf{A} & \textbf{B} & \textbf{AND}
    & \textbf{OR} & \textbf{XOR} & \textbf{NOT A} \\
\midrule
0 & 0 & 0 & 0 & 0 & 1 \\
0 & 1 & 0 & 1 & 1 & 1 \\
1 & 0 & 0 & 1 & 1 & 0 \\
1 & 1 & 1 & 1 & 0 & 0 \\
\bottomrule
\end{tabular}
\end{center}
\end{conceptbox}

%-----------------------------------------------------------------------------
\section{Passing Parameters}
\label{sec:passing-parameters}
%-----------------------------------------------------------------------------

When a function is called, data can be communicated to it through
\textbf{parameters} (also called \textbf{arguments}). C++ provides two
fundamental mechanisms for passing parameters: \emph{by value} and \emph{by
reference}.

%- - - - - - - - - - - - - - - - - - - - - - - - - - - - - - - - - - - - - -
\subsection{Passing by Value}
\label{subsec:pass-by-value}
%- - - - - - - - - - - - - - - - - - - - - - - - - - - - - - - - - - - - - -

\begin{definitionbox}{Pass by Value}
When a parameter is passed \textbf{by value}, a \emph{copy} of the argument's
value is made and given to the function. Any changes the function makes to the
parameter affect only the local copy; the original variable in the calling
function remains unchanged.
\end{definitionbox}

All examples presented so far in this chapter use pass-by-value. Consider the
\texttt{square()} function: when we call \texttt{square(x)}, the value of
\texttt{x} is copied into the parameter \texttt{y}. Modifying \texttt{y}
inside \texttt{square()} would have no effect on \texttt{x} in
\texttt{main()}.

\begin{cppcode}
#include <iostream>
using namespace std;

void tryToChange(int n)
{
    n = 999;   // Only changes the LOCAL copy
    cout << "Inside function: n = " << n << endl;
}

int main()
{
    int x = 5;
    cout << "Before call: x = " << x << endl;

    tryToChange(x);

    cout << "After call:  x = " << x << endl;
    return 0;
}
\end{cppcode}

\noindent\textbf{Output:}
\begin{verbatim}
Before call: x = 5
Inside function: n = 999
After call:  x = 5
\end{verbatim}

\begin{infobox}{When to Use Pass by Value}
Use pass by value when:
\begin{itemize}
    \item The function only needs to \emph{read} the argument.
    \item The argument is a small, inexpensive-to-copy type (e.g.,
          \texttt{int}, \texttt{float}, \texttt{char}).
    \item You want to guarantee that the original variable is not modified.
\end{itemize}
\end{infobox}

%- - - - - - - - - - - - - - - - - - - - - - - - - - - - - - - - - - - - - -
\subsection{Passing by Reference}
\label{subsec:pass-by-reference}
%- - - - - - - - - - - - - - - - - - - - - - - - - - - - - - - - - - - - - -

\begin{definitionbox}{Pass by Reference}
When a parameter is passed \textbf{by reference}, the \emph{address} of the
original variable is passed to the function. The function receives a pointer
(or reference) to the original data, so any changes made inside the function
\emph{directly affect} the original variable.
\end{definitionbox}

Passing by reference is useful when a function needs to modify variables in
the calling scope, or when copying large data structures would be inefficient.

\begin{importantbox}{Advantages of Pass by Reference}
\begin{itemize}
    \item The function can \textbf{modify} the caller's variables directly.
    \item No copy is created, resulting in \textbf{higher execution speed} and
          \textbf{lower memory usage}, especially for large objects.
    \item A function can effectively ``return'' multiple values by modifying
          several referenced parameters.
\end{itemize}
\end{importantbox}

%- - - - - - - - - - - - - - - - - - - - - - - - - - - - - - - - - - - - - -
\subsubsection{Example: Swapping Two Values Using Pointers}
%- - - - - - - - - - - - - - - - - - - - - - - - - - - - - - - - - - - - - -

The classic demonstration of pass-by-reference is a \texttt{swap} function
that exchanges the values of two variables. Using pointers, the function
receives the \emph{addresses} of the variables and modifies them in place.

\begin{cppcode}
#include <iostream>
using namespace std;

void swap(int *a, int *b)
{
    int temp;
    temp = *a;   // Store the value at address a
    *a = *b;     // Copy value at address b into a
    *b = temp;   // Copy saved value into address b
}

int main()
{
    int x = 10, y = 20;

    cout << "Before swap:" << endl;
    cout << "  x = " << x << ", y = " << y << endl;

    swap(&x, &y);   // Pass addresses of x and y

    cout << "After swap:" << endl;
    cout << "  x = " << x << ", y = " << y << endl;

    return 0;
}
\end{cppcode}

\noindent\textbf{Output:}
\begin{verbatim}
Before swap:
  x = 10, y = 20
After swap:
  x = 20, y = 10
\end{verbatim}

\begin{conceptbox}{Tracing the Swap Step by Step}
The following table traces every statement inside \texttt{swap(\&x, \&y)}
where initially \texttt{x\,=\,10} and \texttt{y\,=\,20}:

\begin{center}
\renewcommand{\arraystretch}{1.3}
\begin{tabular}{@{}clccc@{}}
\toprule
\textbf{Step} & \textbf{Statement}
    & \texttt{temp} & \texttt{*a (x)} & \texttt{*b (y)} \\
\midrule
0 & (entry)          & ---  & 10 & 20 \\
1 & \texttt{temp = *a;}  & 10   & 10 & 20 \\
2 & \texttt{*a = *b;}    & 10   & 20 & 20 \\
3 & \texttt{*b = temp;}  & 10   & 20 & 10 \\
\bottomrule
\end{tabular}
\end{center}

\noindent Because the function operates on the \emph{actual memory locations}
of \texttt{x} and \texttt{y}, the changes persist after the function returns.
\end{conceptbox}

\begin{warningbox}{Pass by Value Would Fail Here}
If the swap function used pass-by-value (\texttt{void swap(int a, int b)}),
only the \emph{local copies} would be swapped. The original variables
\texttt{x} and \texttt{y} in \texttt{main()} would remain unchanged.
Always use pass-by-reference (via pointers or C++ references) when a function
must modify its arguments.
\end{warningbox}

%- - - - - - - - - - - - - - - - - - - - - - - - - - - - - - - - - - - - - -
\subsubsection{Alternative: C++ Reference Syntax}
%- - - - - - - - - - - - - - - - - - - - - - - - - - - - - - - - - - - - - -

C++ also supports pass-by-reference using the \texttt{\&} (ampersand)
notation in the parameter list, which provides a cleaner syntax than pointers:

\begin{cppcode}
#include <iostream>
using namespace std;

void swap(int &a, int &b)
{
    int temp = a;
    a = b;
    b = temp;
}

int main()
{
    int x = 10, y = 20;

    cout << "Before swap: x = " << x
         << ", y = " << y << endl;

    swap(x, y);   // No need to pass addresses

    cout << "After swap:  x = " << x
         << ", y = " << y << endl;

    return 0;
}
\end{cppcode}

\noindent\textbf{Output:}
\begin{verbatim}
Before swap: x = 10, y = 20
After swap:  x = 20, y = 10
\end{verbatim}

\begin{infobox}{Pointers vs.\ References}
Both pointers and references achieve pass-by-reference behaviour. C++
references (\texttt{int \&a}) are generally preferred because they are
simpler, safer (cannot be null), and produce cleaner code. Pointers are still
essential in many scenarios, such as dynamic memory allocation, arrays, and
data structures.
\end{infobox}

%-----------------------------------------------------------------------------
\section{Chapter Summary}
\label{sec:ch1-summary}
%-----------------------------------------------------------------------------

\begin{importantbox}{Key Takeaways -- Functions in C++}
\begin{enumerate}
    \item A \textbf{function} groups related statements into a reusable unit,
          promoting modularity, readability, and maintainability.
    \item Every function has a \textbf{return type}, a \textbf{name}, a
          \textbf{parameter list}, and a \textbf{body}.
    \item A \textbf{function prototype} declares a function before its
          definition, allowing flexible source-file organisation.
    \item The \texttt{return} statement sends a value back to the caller
          and terminates the function.
    \item \textbf{Pass by value} copies the argument; the original is safe
          from modification.
    \item \textbf{Pass by reference} (via pointers or \texttt{\&}) gives the
          function direct access to the caller's variable, enabling in-place
          modification.
\end{enumerate}
\end{importantbox}
